\section{Related Work}\label{sec:related}
\paragraph{Multivariate time-series imputation.}
Deep imputation models exploit temporal dependencies and cross-variable relationships through different reconstruction mechanisms. BRITS uses bidirectional recurrent estimation~\citep{brits}, and SAITS employs diagonally masked self-attention~\citep{saits}. CSDI models conditional distributions with score-based diffusion~\citep{csdi}, while ImputeFormer introduces low-rank spatiotemporal attention~\citep{imputeformer}. Our framework integrates contextual estimation with explicit patch evidence and a subsequent observation-guided correction stage.

\paragraph{Temporal and cross-variable representation.}
PatchTST represents local trajectories using temporal patches~\citep{patchtst}; iTransformer models dependencies between variate tokens~\citep{itransformer}. T1 binds temporal feature channels to variable-attention heads~\citep{t1}. PCE adopts this channel--head geometry as its variable-signature attention building block. MPEC and PEGR extend contextual reconstruction with source-organized fragments and pointwise reading, while OGRR revisits the resulting numerical estimate using visible observations.

\paragraph{Structured temporal modeling.}
TimeMixer++ combines multiple representation and mixing operations to capture temporal patterns and analyzes its components through ablation, representation, efficiency and sensitivity studies~\citep{timemixerpp}. For imputation, we organize information around point context, patch evidence and observation-guided correction. Capacity-matched controls, shared refinement and mask-defined error strata connect these design choices to measured reconstruction behavior.
